\documentclass[11pt]{article}
\usepackage[margin=1in]{geometry}
\usepackage{amsmath,amssymb}
\usepackage{enumitem,booktabs}
\usepackage[hidelinks]{hyperref}
\setlist{itemsep=1pt,topsep=2pt}
\newcommand{\lvl}[1]{\textsf{\small[#1]}}
\title{Response to ``Direction for the next phase''\\[2pt]\large Erd\H{o}s Problem \#644}
\author{erdos644 project}
\date{25 November 2026}
\begin{document}
\maketitle

\section*{Summary}
We adopted the minimal-counterexample frame and Statement (G) as the target. Tasks 0, 1, 2, 3 and 5 were carried out
in the order the brief asked; Task 4 is waiting for the papers.
\begin{itemize}
\item \textbf{The lemmas hold.} They passed independent verification and are formalised in Lean.
\item \textbf{The decisive test was negative.} The consequences of saturation do not rule out near-pool families.
\item \textbf{The proof side (Task 2) found no leverage}, within the time limit.
\item \textbf{The disproof side (Task 3) found no construction} with linear excess, and showed why parity-type
constructions cannot give one.
\end{itemize}
The frame gives a cleaner statement of the problem but not, by itself, a way into it. Level tags are as in
\texttt{notes/LEDGER.md}.

\section*{Task 0: verification of \S3}
\begin{itemize}
\item \textbf{Verified} (\texttt{notes/verify/minimal\_frame.md}): Lemmas 1, 3 and 4, Corollary 5, and the equivalence
of (G) with $f(k,7)\le(\tfrac34+o(1))k$. There were no failures; the only gaps were in the proof text: part (ii) of
Lemma 4, $m\ge3$, and $S_B\subseteq V$. The tests covered about 1300 locally minimal families and 1900 random ones,
sampled at $T=2$.
\item \textbf{Added: Lemma 2.5} (a reviewer's point). A minimal counterexample is an antichain. Lemma 4 uses this to
rule out $B\setminus\{z\}$ equalling an existing member.
\item \textbf{Lean.} \texttt{Lean644/MinimalFrame.lean} contains \texttt{private\_set}, \texttt{saturation},
\texttt{antichain} and \texttt{padding\_iff}. The build and the trust check pass, using only the standard axioms.
\item \textbf{Provenance.} Lemma 3 is the classical Bollob\'as set-pair / $\tau$-criticality argument, and the notes now
say so. On its own it forces only $N\ge(1+t)k$. Lemma 4 (saturation) is the new ingredient.
\end{itemize}

\section*{Reviewer refinements adopted during the phase}
\begin{itemize}
\item \textbf{Reading of Task 1.} Looseness is defined only for genuine $(7,2)$-families with $\tau>T$. At
$\varepsilon=0$, shrinking always ends saturated, so Task 1(a) was used as a template rather than a test. For near-pool
and hash-like \emph{models}, we tested whether the consequences of saturation contradict their structure.
\item \textbf{Uniform variant.} Swap $z$ for a fresh point instead of deleting it, using the measure
$(|H|,\sum_x\binom{\deg x}{2})$. The family stays $k$-uniform, which gives Lemmas 3U and 4U. Part (ii) of Lemma 4U
(``no single point covers $B_2,\dots,B_m$'') turned out to be implied by part (i) \lvl{HAND}. So the variant equals
Lemma 4 restricted to points of degree $\ge2$, plus uniformity.
\end{itemize}

\section*{Task 1: are the surviving candidates saturated?}
\begin{itemize}
\item \textbf{(a) FKW, $m=3$.} There are ten incidence classes, and exactly one of them is loose: members with 11 points
of $X$, with $z$ their only $Y$-point \lvl{EXACT}. Shrinking that class gives $H_1$: 323\,236 members, $\tau=10$, an
antichain, saturated in every class. Its $(7,2)$ property is proved by hand and confirmed by MILP. Notably, $H_1$
contains a pencil shape yet has no good seven, so the pencil finish $(\Sigma)$(b) fails for non-uniform families.
\item \textbf{(b) The $k=8$ family.} Its $(7,2)$ status is undecided, so looseness is undefined there. A sample of
incidences found none loose.
\item \textbf{(c) Near-pool models at $(10,8,21)$.} The sets $U_z$ of Corollary 5 exist in 31 of 32 incidence classes
\lvl{EXACT}, and this is unchanged with the uniform condition (ii). Private sets fail in the symmetric models, but
because of the symmetry, not the pool.
\end{itemize}
\textbf{Decision under the brief's rule:} saturation does not make near-pool or hash-like structure incompatible. The
open class does not shrink.

\section*{Task 2: proof side (negative, stopped at the time limit)}
\begin{itemize}
\item \textbf{(a)} No common core of mass about $\omega$ is forced inside $K$. Even a common core would not supply the
pool rule that Theorem P needs; obtaining that is CP itself.
\item \textbf{(b)} The slack summed over all pairs $(K,z)$ is an identity, $\|H\|(4\delta+1)$. A contradiction would
need a lower bound on waste, which saturation does not provide.
\item \textbf{(c)} In the uniform frame, points of degree 1 carry no condition, so no bound on $N$ follows.
\item \textbf{A finite-$k$ obstruction.} The complete families on $\lceil 7k/4\rceil$ points, for $k=5,9,13$, are
saturated in every incidence and satisfy every lemma of the frame \lvl{EXACT}. These are $O(1)$ rounding cases, but
they show that the frame alone cannot separate genuine counterexamples from near-misses.
\end{itemize}

\section*{Task 3: disproof side}
\begin{itemize}
\item \textbf{Template.} $H_1$ made deletion-minimal is $H_1^{\min}$: 100\,506 members, every one with a private set,
$\tau=10$, and $(7,2)$ \lvl{EXACT}. A sample of 1\,800 incidences were all saturated. There is also a uniform template
$U_1$: 12-uniform on 34 points, minimal and saturated (sampled).
\item \textbf{The role of parity} \lvl{HAND}. The parity condition colours the $(k+1)$-layer, so minimality splits into
independent covering problems, one per class. Parity contributes no slack; its only effect is half-point rounding.
\item \textbf{Second constraints} (two parities, products, tensors, conditions mod 3). These are all profile families.
Theorem S$^+$ (verified) then forces a good seven once $T-3k/4\ge 3(q+4)/4$, where $q$ is the number of cells. So their
excess is $O(1)$, never linear. Finite scans agree. No family with $(7,2)$ and $\tau\ge 3k/4+2$ was found.
\end{itemize}

\section*{Task 4: literature}
\begin{itemize}
\item Erd\H{o}s--Hajnal--Tuza (1991) has been read. Its setting counts \emph{points} and treats $s=1$; it does not apply
directly to (7,2).
\item Fon-Der-Flaass--Kostochka--Woodall (1999) has been re-read. Its construction is stated for $k\ge10$; a remark
claims it works for $k\ge4$ and fails for $k=2$, so $f(8,7)\in\{6,7\}$ is left open by the paper. Padding gives
$f(r,7)\ge 3\lfloor r/4\rfloor+1$ for $r\ge40$.
\item EFKT92 and Kostochka (2002) are awaited. When they arrive, we will extract how each proof chooses its extremal
configuration and whether it uses criticality.
\end{itemize}

\section*{Task 5: corrections}
\begin{itemize}
\item The reviewer brief now states the open problem as Statement (G), and presents CP as a necessary special case whose
proof alone does not bound $f(k,7)$.
\item The petal side is described as closed only against the adversary classes analysed.
\item The README, the status report and the external brief were corrected the same way, including the confinement
caveat.
\end{itemize}

\section*{Side result (not counted toward the goal)}
The certification of the one-round sparsest-anchor argument is still running at $t=43/50$. Its continuum cover has more
than 5\,000 regions certified and none failed. The exact closure certificates and finite-$k$ rounding are not finished,
and nothing is claimed yet. If completed, it would give $f(k,7)\le(0.86+o(1))k$ for all families, below the published
$7/8$.

\section*{Assessment}
The phase answered its decisive question in the negative. Private sets and saturation, alone or combined, in the
deletion or uniform form, do not constrain near-pool families. The minimal frame therefore removes the lifting and petal
side problems as separate issues, but it does not supply the missing ingredient for (G). A full resolution still needs a
new idea: either a mechanism that uses $T$-covering adaptively beyond what saturation encodes, or a construction whose
excess over $3k/4$ grows linearly and escapes the profile families covered by Theorem S$^+$.
\end{document}
